\documentclass{article}

\usepackage{amsmath,amssymb}
\usepackage{xurl}
\usepackage[hidelinks]{hyperref}
\setlength{\emergencystretch}{2em}

% Shared commands for notes in docs/paper-gaps/.

\DeclareUnicodeCharacter{2080}{\ifmmode{}_{0}\else\textsubscript{0}\fi}
\DeclareUnicodeCharacter{2081}{\ifmmode{}_{1}\else\textsubscript{1}\fi}
\DeclareUnicodeCharacter{2082}{\ifmmode{}_{2}\else\textsubscript{2}\fi}
\DeclareUnicodeCharacter{2099}{\ifmmode{}_{n}\else\textsubscript{n}\fi}

\newcommand{\Lean}{\textsc{Lean}}
\ExplSyntaxOn
\NewDocumentCommand{\leanid}{m}{
  \ifmmode
    \textsf{\detokenize{#1}}
  \else
    \group_begin:
    \urlstyle{sf}
    \tl_set:Nn \l_tmpa_tl {#1}
    \tl_replace_all:Nnn \l_tmpa_tl {\_} {_}
    \exp_args:NV \path \l_tmpa_tl
    \group_end:
  \fi
}
\ExplSyntaxOff
\providecommand{\tr}{\operatorname{tr}}
\newcommand{\tnleanIssueBase}{https://github.com/LionSR/TNLean/issues/}
\newcommand{\tnleanPrBase}{https://github.com/LionSR/TNLean/pull/}
\newcommand{\tnleanDocBase}{https://sirui-lu.com/TNLean/docs/}
\newcommand{\ghissue}[1]{\href{\tnleanIssueBase#1}{GitHub issue \##1}}
\newcommand{\ghpr}[1]{\href{\tnleanPrBase#1}{PR \##1}}
\newcommand{\doclink}[1]{\href{\tnleanDocBase#1.html}{\path{#1}}}

% Dirac notation and the identity operator, as in blueprint/src/macros/common.tex.
% \providecommand keeps note-local definitions authoritative.
\usepackage{dsfont}
\providecommand{\ket}[1]{|#1\rangle}
\providecommand{\bra}[1]{\langle #1|}
\providecommand{\braket}[2]{\langle #1 | #2 \rangle}
\providecommand{\ketbra}[2]{|#1\rangle\!\langle #2|}
\providecommand{\Id}{\mathds{1}}
\providecommand{\one}{\mathds{1}}
\providecommand{\C}{\mathbb{C}}
\providecommand{\MN}[1]{M_{#1}(\C)}
\providecommand{\MD}{M_D(\C)}

% External referencing for paper sources.  The build helper writes sanitized
% auxiliary files under build/paper-aux/ and excludes duplicated source labels.
\usepackage{xr}

% Import a generated paper auxiliary file with a caller-supplied label prefix.
% The candidate paths support builds from docs/paper-gaps/, the repository root,
% and build/paper-aux/ itself.
\newcommand{\importpaperaux}[2]{%
  \IfFileExists{../../build/paper-aux/#2.aux}{%
    \externaldocument[#1]{../../build/paper-aux/#2}%
  }{%
    \IfFileExists{build/paper-aux/#2.aux}{%
      \externaldocument[#1]{build/paper-aux/#2}%
    }{%
      \IfFileExists{#2.aux}{%
        \externaldocument[#1]{#2}%
      }{}%
    }%
  }%
}

% General external reference macros
\newcommand{\paperref}[2]{\ref{#1-#2}}
\newcommand{\papereqref}[2]{(\ref{#1-#2})}

% CPSV16 source (1606.00608 / MPDO-22-12-17-2.tex) external document and shortcuts
\importpaperaux{cpsv16-}{1606.00608/MPDO-22-12-17-2-tnlean}
\newcommand{\cpsvref}[1]{\paperref{cpsv16}{#1}}
\newcommand{\cpsveqref}[1]{\papereqref{cpsv16}{#1}}

% Verdict marker: \gapnote{<kind>}{<status>}. Kind is the policy
% classification (clarification, local-correction, scope-restriction,
% unfaithful, false-source, open-gap); status is open, wip (elimination
% underway), resolved, or historical. Machine-read by the site index;
% prints a verdict line.
\newcommand{\gapnote}[2]{\par\noindent\textbf{Verdict:} #1 (#2).\par}


\title{Common Identity Coefficients in the Block Ground-Space Intersection}
\author{}
\date{2026-07-30}

\begin{document}
\maketitle
\gapnote{scope-restriction}{resolved}

\section{The source argument}

The proof of Theorem~12 of
Ref.~\cite[lines~1424--1456]{PerezGarcia2007MPSReps} considers canonical
blocks $A^j=(A_a^j)_a$ and a vector in the intersection
\begin{align}
    \left(\C^d\otimes\bigoplus_j\mathcal G_m^{A^j}\right)
    \cap
    \left(\bigoplus_j\mathcal G_m^{A^j}\otimes\C^d\right).
    \label{eq:pgvwc07_common_identity_coefficients_ground_space_intersection}
\end{align}
Here, $\mathcal G_m^{A^j}$ denotes the length-$m$ MPS image space generated by
the block $A^j$. Comparing the two boundary representations of a vector in
(\ref{eq:pgvwc07_common_identity_coefficients_ground_space_intersection})
gives matrices $C_a^j$ and $D_b^j$ satisfying
\begin{align}
    A_b^j C_a^j
        &= D_b^j A_a^j
        \qquad\text{for every block $j$ and all physical indices $a,b$.}
    \label{eq:pgvwc07_common_identity_coefficients_boundary_intertwining}
\end{align}

The source uses the right-canonical normalization
\begin{align}
    \sum_a A_a^j A_a^{j\dagger}
        &= \Id.
    \label{eq:pgvwc07_common_identity_coefficients_right_canonical}
\end{align}
Its displayed definition of $E^j$ omits the adjoint needed by the subsequent
calculation. The adjoint-corrected boundary matrix is
\begin{align}
    E^j
        &= \sum_a C_a^j A_a^{j\dagger}.
    \label{eq:pgvwc07_common_identity_coefficients_corrected_boundary_matrix}
\end{align}
Summing
(\ref{eq:pgvwc07_common_identity_coefficients_boundary_intertwining}) over $a$
after right multiplication by $A_a^{j\dagger}$, and then using
(\ref{eq:pgvwc07_common_identity_coefficients_right_canonical}), gives
\begin{align}
    A_b^j E^j
        &= \sum_a A_b^j C_a^j A_a^{j\dagger},
    \label{eq:pgvwc07_common_identity_coefficients_boundary_sum_left}\\
        &= \sum_a D_b^j A_a^j A_a^{j\dagger},
    \label{eq:pgvwc07_common_identity_coefficients_boundary_sum_intertwined}\\
        &= D_b^j.
    \label{eq:pgvwc07_common_identity_coefficients_boundary_factorization}
\end{align}
Substitution into
(\ref{eq:pgvwc07_common_identity_coefficients_boundary_intertwining}) yields
\begin{align}
    A_b^j C_a^j
        &= A_b^j E^j A_a^j.
    \label{eq:pgvwc07_common_identity_coefficients_source_boundary_identity}
\end{align}
This source-aligned argument is formalized by
\leanid{MPSTensor.pgvwc07_mem_iSup_groundSpace_of_trace_decomposition}.

\section{The common-coefficient variant}

The local variant
\leanid{MPSTensor.pgvwc07_mem_iSup_groundSpace_of_trace_decomposition_of_identity_coefficients}
does not assume the right-canonical normalization
(\ref{eq:pgvwc07_common_identity_coefficients_right_canonical}). Instead, it
assumes scalars $c_a$, independent of the block index $j$, such that
\begin{align}
    \sum_a c_a A_a^j
        &= \Id
        \qquad\text{for every block $j$.}
    \label{eq:pgvwc07_common_identity_coefficients_common_identity_expansion}
\end{align}
For each block, define
\begin{align}
    E^j
        &= \sum_a c_a C_a^j.
    \label{eq:pgvwc07_common_identity_coefficients_variant_boundary_matrix}
\end{align}
Multiplying
(\ref{eq:pgvwc07_common_identity_coefficients_boundary_intertwining}) by $c_a$
and summing over $a$ gives
\begin{align}
    A_b^j E^j
        &= \sum_a c_a A_b^j C_a^j,
    \label{eq:pgvwc07_common_identity_coefficients_weighted_boundary_sum}\\
        &= \sum_a c_a D_b^j A_a^j,
    \label{eq:pgvwc07_common_identity_coefficients_weighted_intertwining}\\
        &= D_b^j.
    \label{eq:pgvwc07_common_identity_coefficients_variant_factorization}
\end{align}
Consequently,
\begin{align}
    A_b^j C_a^j
        &= A_b^j E^j A_a^j.
    \label{eq:pgvwc07_common_identity_coefficients_variant_boundary_identity}
\end{align}
The conclusion
(\ref{eq:pgvwc07_common_identity_coefficients_variant_boundary_identity})
agrees with the local step
(\ref{eq:pgvwc07_common_identity_coefficients_source_boundary_identity}), but
the hypothesis and the construction of $E^j$ differ.

\section{Formal boundary}

The common-coefficient theorem is a conditional local variant, not a
formalization of the right-canonical step in Theorem~12 of
Ref.~\cite{PerezGarcia2007MPSReps}. Results that derive
(\ref{eq:pgvwc07_common_identity_coefficients_common_identity_expansion}) from
a simultaneous one-letter word-span condition may use this variant.

A source-faithful proof should instead use
\leanid{MPSTensor.pgvwc07_mem_iSup_groundSpace_of_trace_decomposition}, which
retains the right-canonical hypothesis and uses the adjoint-corrected boundary
matrix
(\ref{eq:pgvwc07_common_identity_coefficients_corrected_boundary_matrix}).
Neither theorem supplies the hypotheses of the other. Transferring a
hypothesis between them therefore requires a separate implication.

\section{Conclusion}

The common-identity-coefficient result proves the same boundary identity as the
right-canonical source calculation, but under a different local hypothesis.
The distinction is a resolved scope restriction: the formal development
contains both the conditional common-coefficient variant and the
source-aligned right-canonical theorem.

\bibliographystyle{alpha}
\bibliography{references}

\end{document}
